% SOURCE_STATEMENT: "Problem 2. Consider the planar restricted three-body problem. For Jacobi constants such that the Hill's region looks like a barbell, do transit orbits always exist? (Rick Moeckel)"
% Preparation and final self-audit: CONTRIBUTING.md prompts C and D; source/audit artifacts retained in the local physics_statements_2026_10_11 batch.
% !TeX program = lualatex
\documentclass[11pt]{article}
\usepackage[margin=25mm]{geometry}
\usepackage{fontspec}
\setmainfont{Latin Modern Roman}
\usepackage{amsmath,amssymb,mathtools}
\usepackage{unicode-math}
\setmathfont{Latin Modern Math}
\usepackage{xurl,hyperref}
\hypersetup{colorlinks=true,urlcolor=blue}
\setcounter{secnumdepth}{0}
\setlength{\parindent}{0pt}
\setlength{\parskip}{5pt}
\setlength{\emergencystretch}{3em}
\newcommand{\sref}[2]{\hyperlink{source:#1:#2}{[S#2]}}
\newcommand{\cataloguescope}[1]{\paragraph{Recorded target.}#1}
\newcommand{\R}{\mathbb R}
\newcommand{\C}{\mathbb C}
\newcommand{\N}{\mathbb N}
\newcommand{\Z}{\mathbb Z}
\newcommand{\Q}{\mathbb Q}
\newcommand{\F}{\mathbb F}
\newcommand{\eps}{\varepsilon}
\begin{document}
% UnsolvedMath ID 20002531; source catalogue code AIM-PHYSICS-0039
\setcounter{section}{20002530}
\section{A gauge-invariant action-gap certificate for barbell transit orbits}\label{20002531}
{\small UnsolvedMath ID 20002531 · Mathematical Physics}
\subsection{Definitions and mathematical statement}

\paragraph{Shared notation.}$\mathbb N=\{1,2,\ldots\}$; the standard number fields have their usual meanings. All additional conventions are stated locally.


\paragraph{Definitions and notation.}In the circular restricted three-body problem, two primaries follow prescribed circular motion and a test particle has negligible effect on them. In rotating coordinates the Jacobi constant combines kinetic energy and the effective gravitational/centrifugal potential. The Hill region is the set where the zero-velocity inequality permits motion. Its necks occur near collinear equilibria. A barbell-shaped Hill region has two larger permitted components connected by a neck. A transit orbit passes through the connecting region from one binary neighborhood toward another, near a collinear relative equilibrium.

\paragraph{Catalogue wording retained for traceability.}
In the planar restricted three-body problem, do transit orbits always exist for Jacobi constants for which the Hill region has the barbell shape in the source? More generally, identify settings in the $n$-body problem where the Maupertuis--Lagrange variational approach yields new results, accounting for degeneracy at the Hill-region boundary.
\paragraph{Scope and source qualification.}The source records small-neck existence and asks about all relevant Jacobi constants. A variational approach must address the degeneracy of the Jacobi metric at the Hill boundary; a local-neck result is not the full target. The catalogue title/status was paired with a machine-generated research attempt; that attempt is not peer-reviewed evidence of settlement. The statement covers the original source questions rather than only the title’s subcase.
\subsection{Short English statement}
Show whether every barbell Hill region of the planar restricted three-body problem admits transit trajectories.
\subsection{Sources}
\begin{itemize}
\item[\hypertarget{source:20002531:1}{S1}] ULAM.AI and the UnsolvedMath Contributors, UnsolvedMath, record 20002531 (AIM-PHYSICS-0039). Catalogue curation under CC BY 4.0; definitions and English exposition added locally.
\url{https://huggingface.co/datasets/ulamai/UnsolvedMath}.
\item[\hypertarget{source:20002531:2}{S2}] American Institute of Mathematics, Variational Methods in Celestial Mechanics; notes by Kuo-Chang Chen, version 22 June 2003, source item 2.
\url{https://aimath.org/WWN/varcelest/varcelest.pdf}.
\end{itemize}
\label{end:20002531}
\end{document}
